% @see : https://coopmaths.fr/alea/?uuid=a17c6&id=auto6M2C-flash1&n=1&d=10&s=1&s2=false&s3=2&qcm=0&cd=1&alea=Ggek&uuid=b486a&id=auto6M2C-flash2&n=1&d=10&s=1&s2=false&s3=4&s4=2&qcm=0&cd=1&alea=OM6x&uuid=a17d6&id=auto6M2C-flash3&alea=drm2&uuid=83be9&id=auto6M2C-1&n=1&d=10&s=3&s2=false&s3=2&s4=1&cd=1&alea=GAxI&uuid=026d9&id=6M2A&n=1&d=10&s=1&s2=true&cd=1&alea=BOdr&uuid=026d9&id=6M2A&n=1&d=10&s=3&s2=true&cd=1&alea=gNak&uuid=026d9&id=6M2A&n=1&d=10&s=2&s2=true&cd=1&alea=oCF4&uuid=c6fb2&id=6M2C&n=1&d=10&s=1&cd=1&alea=EpJe&uuid=c6fb2&id=6M2C&n=1&d=10&s=2&cd=1&alea=XVGT&uuid=c6fb2&id=6M2C&n=1&d=10&s=3&cd=1&alea=RD7u&uuid=bae27&id=6M2C-flash1&alea=8VtO&uuid=90e79&id=6M2C-flash2&n=1&d=10&cd=1&alea=8N1g&pdfParam=eyJ0aXRsZSI6IiIsInJlZmVyZW5jZSI6IiIsInN1YnRpdGxlIjoiIiwic3R5bGUiOiJDb29wbWF0aHMiLCJmb250T3B0aW9uIjoiU3RhbmRhcmRGb250IiwidGFpbGxlRm9udE9wdGlvbiI6MTIsImR5c1RhaWxsZUZvbnRPcHRpb24iOjE0LCJjb3JyZWN0aW9uT3B0aW9uIjoiQXZlY0NvcnJlY3Rpb24iLCJxcmNvZGVPcHRpb24iOiJTYW5zUXJjb2RlIiwidHlwZUZpY2hlIjoiRmljaGUiLCJkdXJhdGlvbkNhbk9wdGlvbiI6IjkgbWluIiwidGl0bGVPcHRpb24iOiJTYW5zVGl0cmUiLCJuYlZlcnNpb25zIjoxLCJleG9zIjp7fX0\%3D&v=eleve
\documentclass[a4paper,11pt,fleqn]{article}

%%% COULEURS %%%
\usepackage[table,svgnames]{xcolor}
\definecolor{nombres}{cmyk}{0,.8,.95,0}
\definecolor{gestion}{cmyk}{.75,1,.11,.12}
\definecolor{gestionbis}{cmyk}{.75,1,.11,.12}
\definecolor{grandeurs}{cmyk}{.02,.44,1,0}
\definecolor{geo}{cmyk}{.62,.1,0,0}
\definecolor{algo}{cmyk}{.69,.02,.36,0}
\definecolor{correction}{cmyk}{.63,.23,.93,.06}
\arrayrulecolor{couleur_theme} % Couleur des filets des tableaux

\usepackage[most]{tcolorbox} % Supprimé à cause de conflits avec ProfCollege
\tcbuselibrary{documentation}

%%%%%%%% Modification paramétrage pour footnotes
% Le paquet semble déjà appelé en amont, curieux qu'il n'y ait pas eu de clash
% avec l'appel ci-dessous mais seulement quand j'ai essayé de l'inclure avec l'option lualatex

% \usepackage{hyperref}
% Parametrages
\hypersetup{
    colorlinks=true,% On active la couleur pour les liens. Couleur par défaut rouge
    linkcolor=black,% On définit la couleur pour les liens internes
    % filecolor=magenta,% On définit la couleur pour les liens vers les fichiers locaux      
    urlcolor=black,% On définit la couleur pour les liens vers des sites web
    % pdftitle={Puissance Quatre},% On définit un titre pour le document pdf
    % pdfpagemode=FullScreen,% On fixe l'affichage par défaut à plein écran
    }
% Pour avoir les numérotations arabic y compris dans les environnements tcolorbox
\renewcommand{\thempfootnote}{\arabic{mpfootnote}}
%%%%%%%%

\usepackage{multicol} 
\usepackage{calc}
\usepackage{enumitem}
\usepackage{graphicx}
\usepackage{tabularx}
\usepackage{tablvar}

\setlength{\parindent}{0mm}
\renewcommand{\arraystretch}{1.5}
\renewcommand{\labelenumi}{\textbf{\theenumi{}.}}
\renewcommand{\labelenumii}{\textbf{\theenumii{}.}}
\setlength{\fboxsep}{3mm}

\setlength{\headheight}{14.5pt}

\setlength{\premulticols}{0pt}

\newcounter{ExoMA}

\newlength{\parindentMA}%
\setlength{\parindentMA}{\parindent}
\addtolength{\parindentMA}{30pt}

\usepackage{simplekv}

\setKVdefault[Theme]{Coopmath,Classiques=false,CANs=false}
\defKV[Theme]{Classique=\setKV[Theme]{Coopmath=false}\setKV[Theme]{Classiques}}
\defKV[Theme]{CAN=\setKV[Theme]{Coopmath=false}\setKV[Theme]{CANs}}

\usepackage{fancyhdr}
\fancyhead[C]{}
\fancyfoot{}

\def\bla{}

\NewDocumentCommand\Theme{o m m m m}{%
  \useKVdefault[Theme]%
  \setKV[Theme]{#1}%
  \ifboolKV[Theme]{CANs}{%
    \usepackage[left=1.5cm,right=1.5cm,top=2cm,bottom=2cm]{geometry}%
    \fancypagestyle{premierePage}{%
      \fancyhead[C]{\textsc{\ifx\bla#2\bla Course aux nombres\else#2\fi}}%
      \fancyfoot[R]{\scriptsize Coopmaths.fr -- CC-BY-SA}%
      \setlength{\headheight}{14.5pt}%
      \NewDocumentCommand\dureeCan{}{\ifx\bla#4\bla 9 minutes\else #4\fi}
      \NewDocumentCommand\titreSujetCan{}{\ifx\bla#3\bla\else\textbf{#3}\fi}
    }%
    \pagestyle{premierePage}
    \colorlet{couleur_theme}{black}%
    \colorlet{couleur_numerotation}{couleur_theme}%
    \let\EXO\EXOCan\let\endEXO\endEXOCan%
  }{%
    \renewcommand{\headrulewidth}{0pt}
    \renewcommand{\footrulewidth}{0pt}
    \pagestyle{fancy}
    \ifboolKV[Theme]{Classiques}{%
      \usepackage[left=0.65cm,right=0.65cm,top=2cm,bottom=1.5cm]{geometry}
      \let\EXO\EXOlibre\let\endEXO\endEXOlibre%
      \colorlet{couleur_theme}{black}%
      \colorlet{couleur_numerotation}{couleur_theme}%
      \fancyhead[C]{\bfseries #3}
      \fancyhead[L]{\bfseries #4}
      \fancyfoot[R]{#5}
    }{%
      \usepackage[left=1.5cm,right=1.5cm,top=4cm,bottom=2cm]{geometry}
      \theme{#2}{#3}{#4}{#5}%
      \let\EXO\EXOcoop\let\endEXO\endEXOcoop%
    }%
  }%
}%

\NewDocumentEnvironment{EXOCan}{m m +b}{%
  #3
}

\NewDocumentEnvironment{EXOcoop}{m m +b}{%
  \begin{tcolorbox}[%
    enhanced,
    breakable,
    colback=white,
    colframe=white,
    coltitle=black,
    title=\IfNoValueTF{#1}{}{\textmd{#1}},%
    overlay unbroken and first={%
      \IfNoValueTF{#2}{}{%
        \node[%
        fill=white,
        anchor=east,
        xshift=10pt,
        text=black
        ]
        at (frame.north east)
        {\scriptsize #2};
      }
      \node[anchor=west,xshift=-20pt,yshift=-15pt] at (frame.north west) {\stepcounter{ExoMA}\numb{\arabic{ExoMA}}};
    }
    ]
    #3
  \end{tcolorbox}
}

\NewDocumentEnvironment{EXOlibre}{m m +b}{%
  \begin{tcolorbox}[%
    enhanced,
    breakable,
    colback=white,
    colframe=white,%couleur_theme,
    coltitle=black,
    title=\stepcounter{ExoMA}\textbf{Exercice \arabic{ExoMA}},%
    ]
    \IfNoValueTF{#1}{}{#1\par}%
    #3
  \end{tcolorbox}%
}


%%% MISE EN PAGE %%%
\usepackage{setspace}
\usepackage{pgf}%
\usepackage{pgfplots}
\pgfplotsset{compat=1.18}
\usetikzlibrary{babel,arrows, arrows.meta,calc,fit,patterns,plotmarks,shapes.geometric,shapes.misc,shapes.symbols,shapes.arrows,shapes.callouts, shapes.multipart, shapes.gates.logic.US,shapes.gates.logic.IEC, er, automata,backgrounds,chains,topaths,trees,petri,mindmap,matrix, calendar, folding,fadings,through,positioning,scopes,decorations.fractals,decorations.shapes,decorations.text,decorations.pathmorphing,decorations.pathreplacing,decorations.footprints,decorations.markings,shadows}

\spaceskip=2\fontdimen2\font plus 3\fontdimen3\font minus3\fontdimen4\font\relax %Pour doubler l'espace entre les mots
\newcommand\numb[1]{ % Dessin autour du numéro d'exercice
  \begin{tikzpicture}[overlay,scale=.8]%,yshift=-.3cm
    \draw[fill=couleur_numerotation,couleur_numerotation](-.4,-0.1)rectangle($(-0.4,.8)+(2em,0)$);
    \draw (-.4,.8) node[white,anchor=north west,inner sep=0.5pt]{\bfseries EX}; 
    \draw[line width=.05cm,couleur_numerotation,fill=white] (0,0)--(.5,.5)--(1,0)--(.5,-.5)--cycle;
    \draw (0.5,0) node[anchor=center,couleur_numerotation]{\large\bfseries#1};
  \end{tikzpicture}
}

% echelle pour le dé
\def\globalscale {0.04}
% abscisse initiale pour les chevrons
\def\xini {3}

\newcommand\theme[4]{%
  \fancyhead[C]{%
    % Tracé du dé
    \begin{tikzpicture}[y=0.80pt, x=0.80pt, yscale=-\globalscale, xscale=\globalscale,remember picture, overlay,transform canvas={xshift=-0.5\linewidth,yshift=2.7cm},fill=couleur_theme]%,xshift=17cm,yshift=9.5cm
      %%%% Arc supérieur gauche%%%%
      \path[fill](523,1424)..controls(474,1413)and(404,1372)..(362,1333)..controls(322,1295)and(313,1272)..(331,1254)..controls(348,1236)and(369,1245)..(410,1283)..controls(458,1328)and(517,1356)..(575,1362)..controls(635,1368)and(646,1375)..(643,1404)..controls(641,1428)and(641,1428)..(596,1430)..controls(571,1431)and(538,1428)..(523,1424)--cycle;
      %%%% Dé face supérieur%%%%
      \path[fill](512,1272)..controls(490,1260)and(195,878)..(195,861)..controls(195,854)and(198,846)..(202,843)..controls(210,838)and(677,772)..(707,772)..controls(720,772)and(737,781)..(753,796)..controls(792,833)and(1057,1179)..(1057,1193)..controls(1057,1200)and(1053,1209)..(1048,1212)..controls(1038,1220)and(590,1283)..(551,1282)..controls(539,1282)and(521,1278)..(512,1272)--cycle;
      %%%% Dé faces gauche et droite%%%%
      \path[fill](1061,1167)..controls(1050,1158)and(978,1068)..(900,967)..controls(792,829)and(756,777)..(753,756)--(748,729)--(724,745)..controls(704,759)and(660,767)..(456,794)..controls(322,813)and(207,825)..(200,822)..controls(193,820)and(187,812)..(187,804)..controls(188,797)and(229,688)..(279,563)..controls(349,390)and(376,331)..(391,320)..controls(406,309)and(462,299)..(649,273)..controls(780,254)and(897,240)..(907,241)..controls(918,243)and(927,249)..(928,256)..controls(930,264)and(912,315)..(889,372)..controls(866,429)and(848,476)..(849,477)..controls(851,479)and(872,432)..(897,373)..controls(936,276)and(942,266)..(960,266)..controls(975,266)and(999,292)..(1089,408)..controls(1281,654)and(1290,666)..(1290,691)..controls(1290,720)and(1104,1175)..(1090,1180)..controls(1085,1182)and (1071,1176)..(1061,1167)--cycle;
      %%%% Arc inférieur bas%%%%
      \path[fill](1329,861)..controls(1316,848)and(1317,844)..(1339,788)..controls(1364,726)and(1367,654)..(1347,591)..controls(1330,539)and(1338,522)..(1375,526)..controls(1395,528)and(1400,533)..(1412,566)..controls(1432,624)and(1426,760)..(1401,821)..controls(1386,861)and(1380,866)..(1361,868)..controls(1348,870)and(1334,866)..(1329,861)--cycle;
      %%%% Arc inférieur gauche%%%%
      \path[fill](196,373)..controls(181,358)and(186,335)..(213,294)..controls(252,237)and(304,190)..(363,161)..controls(435,124)and(472,127)..(472,170)..controls(472,183)and(462,192)..(414,213)..controls(350,243)and(303,283)..(264,343)..controls(239,383)and(216,393)..(196,373)--cycle;
    \end{tikzpicture}
    \begin{tikzpicture}[line cap=round,line join=round,remember picture, overlay, shift={(current page.north west)},yshift=-8.5cm]
      \fill[fill=couleur_theme] (0,5) rectangle (21,6);
      \fill[fill=couleur_theme] (\xini,6)--(\xini+1.5,6)--(\xini+2.5,7)--(\xini+1.5,8)--(\xini,8)--(\xini+1,7)-- cycle;
      \fill[fill=couleur_theme] (\xini+2,6)--(\xini+2.5,6)--(\xini+3.5,7)--(\xini+2.5,8)--(\xini+2,8)--(\xini+3,7)-- cycle;  
      \fill[fill=couleur_theme] (\xini+3,6)--(\xini+3.5,6)--(\xini+4.5,7)--(\xini+3.5,8)--(\xini+3,8)--(\xini+4,7)-- cycle;   
      \node[color=white] at (10.5,5.5) {\LARGE \bfseries{ \MakeUppercase{#4}}};
    \end{tikzpicture}
    \begin{tikzpicture}[remember picture,overlay]
      \node[anchor=north east,inner sep=0pt] at ($(current page.north east)+(0,-.8cm)$) {};
      \node[anchor=west, fill=white, inner sep = 0pt] at ($(current page.north west)+(0.2,-2.3cm)$) {\footnotesize \bfseries{MathALÉA}};
    \end{tikzpicture}
    \begin{tikzpicture}[remember picture,overlay]
      \node[anchor=north east,inner sep=0pt] at ($(current page.north east)+(0,-.8cm)$) {};
      \node[anchor=east, fill=white] at ($(current page.north east)+(-2,-1.5cm)$) {\Huge \textcolor{couleur_theme}{\bfseries{\#}} \bfseries{#2} \textcolor{couleur_theme}{\bfseries \MakeUppercase{#3}}};
    \end{tikzpicture}
  }%
  \fancyfoot[R]{%
    \begin{tikzpicture}[remember picture,overlay]
      \node[anchor=south east] at ($(current page.south east)+(-2,0.25cm)$) {\scriptsize {\bfseries \href{https://coopmaths.fr/}{Coopmaths.fr} -- \href{http://creativecommons.fr/licences/}{CC-BY-SA}}};
    \end{tikzpicture}
    \begin{tikzpicture}[line cap=round,line join=round,remember picture, overlay,xscale=0.5,yscale=0.5, shift={(current page.south west)},xshift=35.7cm,yshift=-6cm]
      \fill[fill=couleur_theme] (\xini,6)--(\xini+1.5,6)--(\xini+2.5,7)--(\xini+1.5,8)--(\xini,8)--(\xini+1,7)-- cycle;
      \fill[fill=couleur_theme] (\xini+2,6)--(\xini+2.5,6)--(\xini+3.5,7)--(\xini+2.5,8)--(\xini+2,8)--(\xini+3,7)-- cycle;  
      \fill[fill=couleur_theme] (\xini+3,6)--(\xini+3.5,6)--(\xini+4.5,7)--(\xini+3.5,8)--(\xini+3,8)--(\xini+4,7)-- cycle;  
    \end{tikzpicture}
  }
  \fancyfoot[C]{}
  \colorlet{couleur_theme}{#1}
  \colorlet{couleur_numerotation}{couleur_theme}
  \def\iconeobjectif{icone-objectif-#1}
  \def\urliconeomethode{icone-methode-#1}
}%

\newcommand*\tikzfootMA{%
  \ifboolKV[Theme]{CANs}{
    \begin{tikzpicture}[remember picture,overlay,shift=(current page.north west)]
      \begin{scope}[x=\textwidth-\oddsidemargin,y=\textheight+5pt]
        \draw[correction,line width=4pt,dashed,dash pattern= on 10pt off 10pt,shift={(-\oddsidemargin,\topmargin)}] (0,0) rectangle (1,-1);
      \end{scope}
    \end{tikzpicture} 
  }{%
  \ifboolKV[Theme]{Classiques}{%
  \begin{tikzpicture}[remember picture,overlay,shift=(current page.south west)]
    \begin{scope}[x={(current page.south east)},y={(current page.north west)}]
      \draw[correction,line width=4pt,dashed,dash pattern= on 10pt off 10pt] ($(0,0)+(5mm,1cm)$) rectangle ($(1,1)+(-5mm,-1.75cm)$);
    \end{scope}
  \end{tikzpicture} 
  }{%
  \begin{tikzpicture}[remember picture,overlay,shift=(current page.south west)]
    \begin{scope}[x={(current page.south east)},y={(current page.north west)}]
      \draw[correction,line width=4pt,dashed,dash pattern= on 10pt off 10pt] ($(0,0)+(5mm,1.5cm)$) rectangle ($(1,1)+(-5mm,-3.75cm)$);
    \end{scope}
  \end{tikzpicture} 
  }%
  }
}

\newenvironment{Correction}{%
  \setcounter{ExoMA}{0}%
  \cfoot{\tikzfootMA}
}{}%

\newcommand\version[1]{
  \fancyhead[R]{
    \begin{tikzpicture}[remember picture,overlay]
    \node[anchor=north east,inner sep=0pt] at ($(current page.north east)+(-.5,-.5cm)$) {\large \textcolor{couleur_theme}{\bfseries V#1}};
    \end{tikzpicture}
  }
}

\usepackage[french]{babel}
% loadPackagesFromContent
\usepackage{tikz}
\newcommand{\e}{\mathrm{e}}
\usepackage{amsmath}
\usepackage{amssymb}
\usetikzlibrary{patterns,patterns.meta}
\usepackage{etoolbox}
\newbool{dys}
\setbool{dys}{false}          
\ifbool{dys}{
% POLICE DYS
% ===== VARIABLE =====
\def\FontChoisie{Fira} % valeurs possibles : Fira, lmodern, tgheros
\newcommand{\choiceFontsDys}[1]{
\ifstrequal{#1}{Fira}{%
  % Fira Sans + Fira Math
  % Description : Fira Sans est une police moderne, et Fira Math est une version mathématique compatible qui maintient le style sans-serif.
  % Utilisation : Fonctionne bien avec XeLaTeX et LuaLaTeX.
  \usepackage{fontspec}
  \setmainfont{Fira Sans}
  \setsansfont{Fira Sans}
  \usepackage{unicode-math}
  \setmathfont{Fira Math}
}{}
\ifstrequal{#1}{lmodern}{%
  % Latin Modern Sans
  % Description : Une version sans-serif de la célèbre police Latin Modern, qui est compatible avec mathastext.
  % Utilisation : Fonctionne avec pdfLaTeX, XeLaTeX et LuaLaTeX.
  \usepackage{lmodern}
  \renewcommand{\familydefault}{\sfdefault}
  \usepackage{mathastext}
}{}
\ifstrequal{#1}{tgheros}{%
  % TeX Gyre Heros + mathastext
  % Description : TeX Gyre Heros est une alternative à Helvetica, et le package mathastext permet d'utiliser la même police pour les mathématiques.
  % Utilisation : Fonctionne avec pdfLaTeX, XeLaTeX, ou LuaLaTeX.
  \usepackage{tgheros}
  \renewcommand{\familydefault}{\sfdefault}
  \usepackage{mathastext}
}{}
}
% Fira ou lmodern ou tgheros
\choiceFontsDys{\FontChoisie}

%\usepackage{unicode-math}
%\usepackage{fontspec}
%\setmainfont{TeX Gyre Schola}
%\setsansfont{TeX Gyre Schola}
%\setmathfont{TeX Gyre Schola Math}
%\setmainfont{OpenDyslexic}[Scale=1.0]
%\setsansfont{OpenDyslexic}[Scale=1.0]
%\setmathfont{OpenDyslexic}[Scale=1.0,range=up/{Latin,latin,num}]
\usepackage[fontsize=14]{scrextend}
\usepackage{setspace}
\setstretch{1.7}
% Une valeur d'environ 1.2 à 1.7 est souvent recommandée. Cela permet d'augmenter l'espace entre les lignes tout en offrant un léger espacement entre les mots.
\setlength{\spaceskip}{1.2em}  
% Une valeur d'environ 1.2em à 1.5em est couramment conseillée. Cela crée un espace plus ample entre les mots, ce qui peut aider à réduire la fatigue visuelle et à améliorer la fluidité de la lecture.
}{
% POLICE STANDARD
\usepackage[T1]{fontenc} 
\usepackage[scaled=1]{helvet}
\usepackage[fontsize=12]{scrextend}
}

\Theme[Coopmaths]{nombres}{}{}{}
\begin{document}

% @see : https://coopmaths.fr/alea?uuid=a17c6&id=auto6M2C-flash1&n=1&d=10&s=1&s2=false&s3=2&alea=Ggek&cd=1&cols=1
\begin{EXO}{}{auto6M2C-flash1}

 \begin{tikzpicture}[baseline,scale = 0.5]

    \tikzset{
      point/.style={
        thick,
        draw,
        cross out,
        inner sep=0pt,
        minimum width=5pt,
        minimum height=5pt,
      },
    }
    \clip (-0.1,-0.1) rectangle (8.85,4.1);
    	\draw[color={red},line width = 2,preaction={fill,color = {orange}, opacity = 0.5}] (0,1)--(0,2)--(2,2)--(2,4)--(3,4)--(3,1)--(6,1)--(6,0)--(2,0)--(2,1)--cycle;
	\draw[color={black},preaction={fill,color = {gray}, opacity = 0.5}] (7,3)--(8,3)--(8,4)--(7,4)--cycle;
	\draw [color={black}] (7.5,2.2) node[anchor = center,scale=1, rotate = 0] {1 cm²};

\end{tikzpicture}\\Quelle est l'aire de la figure ci-dessus ?
\end{EXO}

% @see : https://coopmaths.fr/alea?uuid=b486a&id=auto6M2C-flash2&n=1&d=10&s=1&s2=false&s3=4&s4=2&alea=OM6x&cd=1&cols=1
\begin{EXO}{}{auto6M2C-flash2}

 \begin{tikzpicture}[baseline,scale = 0.5]

    \tikzset{
      point/.style={
        thick,
        draw,
        cross out,
        inner sep=0pt,
        minimum width=5pt,
        minimum height=5pt,
      },
    }
    \clip (-0.10000000000000045,-0.1) rectangle (6.85,4.1);
    	\draw[color ={gray},opacity = 0.6] (0,0)--(0,5);
	\draw[color ={gray},opacity = 0.6] (1,0)--(1,5);
	\draw[color ={gray},opacity = 0.6] (2,0)--(2,5);
	\draw[color ={gray},opacity = 0.6] (3,0)--(3,5);
	\draw[color ={gray},opacity = 0.6] (4,0)--(4,5);
	\draw[color ={gray},opacity = 0.6] (5,0)--(5,5);
	\draw[color ={gray},opacity = 0.6] (6,0)--(6,5);
	\draw[color ={gray},opacity = 0.6] (7,0)--(7,5);
	\draw[color ={gray},opacity = 0.6] (8,0)--(8,5);
	\draw[color ={gray},opacity = 0.6] (0,0)--(8,0);
	\draw[color ={gray},opacity = 0.6] (0,1)--(8,1);
	\draw[color ={gray},opacity = 0.6] (0,2)--(8,2);
	\draw[color ={gray},opacity = 0.6] (0,3)--(8,3);
	\draw[color ={gray},opacity = 0.6] (0,4)--(8,4);
	\draw[color ={gray},opacity = 0.6] (0,5)--(8,5);
	\draw[color={black},preaction={fill,color = {gray}, opacity = 0.5}] (6,3)--(6.5,3)--(6.5,3.5)--(6,3.5)--cycle;
	\draw[opacity=1] (6.5,2.5) node[anchor = center, rotate=0] {\scriptsize \color{black}$1 \text{cm}^2$};
	\draw[color={black},preaction={fill,color = {orange}, opacity = 0.5}] (0,0)--(4,0)--(4,4)--(0,4)--cycle;

\end{tikzpicture}\\
 Quelle est l'aire de la figure ci-dessus ?
\end{EXO}

% @see : https://coopmaths.fr/alea?uuid=a17d6&id=auto6M2C-flash3&n=1&d=10&s=2&s2=true&alea=drm2&cd=1&cols=1
\begin{EXO}{}{auto6M2C-flash3}

 L'aire du rectangle est de $35$ cm².\\
            \begin{tikzpicture}[baseline,scale = 0.5]

    \tikzset{
      point/.style={
        thick,
        draw,
        cross out,
        inner sep=0pt,
        minimum width=5pt,
        minimum height=5pt,
      },
    }
    \clip (-0.10000000000000012,-0.1) rectangle (9.85,5.1);
    	\draw[color={black},preaction={fill,color = {white}},pattern color = {gray}, pattern = {Lines[angle=45, distance=5pt, line width=0.3pt]}] (0,0)--(7,0)--(7,5)--(0,5)--cycle;
	\draw[color={black},preaction={fill,color = {orange}}] (3,3)--(4,3)--(4,4)--(3,4)--cycle;
	\draw[color={black},preaction={fill,color = {orange}}] (3,2)--(4,2)--(4,3)--(3,3)--cycle;
	\draw[color={black},preaction={fill,color = {orange}}] (2,2)--(3,2)--(3,3)--(2,3)--cycle;
	\draw[color={black},preaction={fill,color = {orange}}] (4,2)--(5,2)--(5,3)--(4,3)--cycle;
	\draw[color={black},preaction={fill,color = {orange}}] (1,2)--(2,2)--(2,3)--(1,3)--cycle;
	\draw[color={black},preaction={fill,color = {orange}}] (4,3)--(5,3)--(5,4)--(4,4)--cycle;
	\draw[color={black},preaction={fill,color = {orange}}] (0,2)--(1,2)--(1,3)--(0,3)--cycle;
	\draw[color={black},preaction={fill,color = {orange}}] (1,1)--(2,1)--(2,2)--(1,2)--cycle;
	\draw[color={black},preaction={fill,color = {orange}}] (3,1)--(4,1)--(4,2)--(3,2)--cycle;
	\draw[color={black},preaction={fill,color = {orange}}] (2,1)--(3,1)--(3,2)--(2,2)--cycle;
	\draw[color={black},preaction={fill,color = {orange}}] (1,0)--(2,0)--(2,1)--(1,1)--cycle;
	\draw[color={black},preaction={fill,color = {orange}}] (5,3)--(6,3)--(6,4)--(5,4)--cycle;
	\draw[color={black},preaction={fill,color = {orange}}] (6,3)--(7,3)--(7,4)--(6,4)--cycle;
	\draw[color={black},preaction={fill,color = {orange}}] (5,4)--(6,4)--(6,5)--(5,5)--cycle;
	\draw[color={black},preaction={fill,color = {gray}, opacity = 0.5}] (8,4)--(9,4)--(9,5)--(8,5)--cycle;
	\draw [color={black}] (8.5,3.2) node[anchor = center,scale=1, rotate = 0] {1 cm²};

\end{tikzpicture}\\
       Quelle est l'aire de la zone hachurée ?
      
\end{EXO}

% @see : https://coopmaths.fr/alea?uuid=83be9&id=auto6M2C-1&n=1&d=10&s=3&s2=false&s3=2&s4=1&alea=GAxI&cd=1&cols=1
\begin{EXO}{}{auto6M2C-1}

 \begin{tikzpicture}[baseline,scale = 0.5]

    \tikzset{
      point/.style={
        thick,
        draw,
        cross out,
        inner sep=0pt,
        minimum width=5pt,
        minimum height=5pt,
      },
    }
    \clip (-1,-1) rectangle (18,8);
    	\draw[color ={gray},opacity = 0.6] (-1,-1)--(-1,8);
	\draw[color ={gray},opacity = 0.6] (0,-1)--(0,8);
	\draw[color ={gray},opacity = 0.6] (1,-1)--(1,8);
	\draw[color ={gray},opacity = 0.6] (2,-1)--(2,8);
	\draw[color ={gray},opacity = 0.6] (3,-1)--(3,8);
	\draw[color ={gray},opacity = 0.6] (4,-1)--(4,8);
	\draw[color ={gray},opacity = 0.6] (5,-1)--(5,8);
	\draw[color ={gray},opacity = 0.6] (6,-1)--(6,8);
	\draw[color ={gray},opacity = 0.6] (7,-1)--(7,8);
	\draw[color ={gray},opacity = 0.6] (8,-1)--(8,8);
	\draw[color ={gray},opacity = 0.6] (9,-1)--(9,8);
	\draw[color ={gray},opacity = 0.6] (10,-1)--(10,8);
	\draw[color ={gray},opacity = 0.6] (11,-1)--(11,8);
	\draw[color ={gray},opacity = 0.6] (12,-1)--(12,8);
	\draw[color ={gray},opacity = 0.6] (13,-1)--(13,8);
	\draw[color ={gray},opacity = 0.6] (14,-1)--(14,8);
	\draw[color ={gray},opacity = 0.6] (15,-1)--(15,8);
	\draw[color ={gray},opacity = 0.6] (16,-1)--(16,8);
	\draw[color ={gray},opacity = 0.6] (17,-1)--(17,8);
	\draw[color ={gray},opacity = 0.6] (18,-1)--(18,8);
	\draw[color ={gray},opacity = 0.6] (-1,-1)--(18,-1);
	\draw[color ={gray},opacity = 0.6] (-1,0)--(18,0);
	\draw[color ={gray},opacity = 0.6] (-1,1)--(18,1);
	\draw[color ={gray},opacity = 0.6] (-1,2)--(18,2);
	\draw[color ={gray},opacity = 0.6] (-1,3)--(18,3);
	\draw[color ={gray},opacity = 0.6] (-1,4)--(18,4);
	\draw[color ={gray},opacity = 0.6] (-1,5)--(18,5);
	\draw[color ={gray},opacity = 0.6] (-1,6)--(18,6);
	\draw[color ={gray},opacity = 0.6] (-1,7)--(18,7);
	\draw[color ={gray},opacity = 0.6] (-1,8)--(18,8);
	\draw[color={black},preaction={fill,color = {orange}, opacity = 0.5}] (0,0)--(6,0)--(0,3)--cycle;
	\draw[color={black},preaction={fill,color = {orange}, opacity = 0.5}] (7,0)--(14,0)--(14,4)--cycle;
	\draw [color={black}] (3.5,-0.5) node[anchor = center,scale=1, rotate = 0] {Figure 1};
	\draw [color={black}] (10.5,-0.5) node[anchor = center,scale=1, rotate = 0] {Figure 2};
	\draw[color={black},preaction={fill,color = {gray}, opacity = 0.5}] (16,6)--(17,6)--(17,6.5)--(16,6.5)--cycle;
	\draw [color={black}] (16.5,5.5) node[anchor = center,scale=1, rotate = 0] {1 u.a};

\end{tikzpicture}\\\medskip

      Calculer l'aire de la figure 1 et l'aire de la figure 2 en écrivant les calculs.
\end{EXO}

% @see : https://coopmaths.fr/alea?uuid=026d9&id=6M2A&n=1&d=10&s=1&s2=true&alea=BOdr&cd=1&cols=1
\begin{EXO}{Convertir.}{6M2A}

 $0{,}05 \text{ m}^2 = \ldots\ldots\ldots\ldots ~\text{dm}^2$
\end{EXO}

% @see : https://coopmaths.fr/alea?uuid=026d9&id=6M2A&n=1&d=10&s=3&s2=true&alea=gNak&cd=1&cols=1
\begin{EXO}{Convertir.}{6M2A}

 $0{,}02 ~\text{dm}^2 = \ldots\ldots\ldots\ldots \text{ cm}^2$
\end{EXO}

% @see : https://coopmaths.fr/alea?uuid=026d9&id=6M2A&n=1&d=10&s=2&s2=true&alea=oCF4&cd=1&cols=1
\begin{EXO}{Convertir.}{6M2A}

 $0{,}04 ~\text{dm}^2 = \ldots\ldots\ldots\ldots \text{ m}^2$
\end{EXO}

% @see : https://coopmaths.fr/alea?uuid=c6fb2&id=6M2C&n=1&d=10&s=1&alea=EpJe&cd=1&cols=1
\begin{EXO}{\begin{tikzpicture}[baseline,scale = 0.75]

    \tikzset{
      point/.style={
        thick,
        draw,
        cross out,
        inner sep=0pt,
        minimum width=5pt,
        minimum height=5pt,
      },
    }
    \clip (-2,-3) rectangle (22,7);
    	\draw[color={black}] (0,0)--(3.978088,-0.418114)--(4.396201,3.559974)--(0.418114,3.978088)--cycle;
	\draw [color={black}] (-0.389795,-0.31565) node[anchor = center,scale=1, rotate = 0] {F};
	\draw [color={black}] (4.293738,-0.807909) node[anchor = center,scale=1, rotate = 0] {G};
	\draw [color={black}] (4.785997,3.875624) node[anchor = center,scale=1, rotate = 0] {H};
	\draw [color={black}] (0.102464,4.367883) node[anchor = center,scale=1, rotate = 0] {I};
	\draw[color={black},line width = 0.5] (3.978088,-0.418114)--(3.580279,-0.376302)--(3.62209,0.021506)--(4.019899,-0.020305)--cycle;
	\draw[color={black},line width = 0.5] (0.418114,3.978088)--(0.376302,3.580279)--(0.774111,3.538467)--(0.815923,3.936276)--cycle;
	\draw[color={black},line width = 0.5] (4.396201,3.559974)--(3.998393,3.601785)--(3.956581,3.203976)--(4.35439,3.162165)--cycle;
	\draw[color={black},line width = 0.5] (0,0)--(0.397809,-0.041811)--(0.43962,0.355997)--(0.041811,0.397809)--cycle;
	\draw [color={blue}] (1.989044,-0.209057) node[anchor = center,scale=1, rotate = -6] {//};
\draw [color={blue}] (4.187145,1.57093) node[anchor = center,scale=1, rotate = 84] {//};
\draw [color={blue}] (2.407158,3.769031) node[anchor = center,scale=1, rotate = -6] {//};
\draw [color={blue}] (0.209057,1.989044) node[anchor = center,scale=0.5, rotate = -276] {//};

	\draw [color={black}] (1.93678,-0.706318) node[anchor = center,scale=1, rotate = -6.023989999999998] {4 cm};

\end{tikzpicture}\\
}{6M2C}

 Calculer l'aire du carré.
\end{EXO}

% @see : https://coopmaths.fr/alea?uuid=c6fb2&id=6M2C&n=1&d=10&s=2&alea=XVGT&cd=1&cols=1
\begin{EXO}{\begin{tikzpicture}[baseline,scale = 0.75]

    \tikzset{
      point/.style={
        thick,
        draw,
        cross out,
        inner sep=0pt,
        minimum width=5pt,
        minimum height=5pt,
      },
    }
    \clip (-2,-3) rectangle (22,7);
    	\draw[color={black}] (8,0)--(9.980536,-0.278346)--(10.676402,4.672994)--(8.695866,4.95134)--cycle;
	\draw [color={black}] (7.752786,-0.431636) node[anchor = center,scale=1, rotate = 0] {J};
	\draw [color={black}] (10.099199,-0.761403) node[anchor = center,scale=1, rotate = 0] {K};
	\draw [color={black}] (10.923616,5.10463) node[anchor = center,scale=1, rotate = 0] {L};
	\draw [color={black}] (8.577203,5.434397) node[anchor = center,scale=1, rotate = 0] {M};
	\draw[color={black},line width = 0.5] (9.980536,-0.278346)--(9.584429,-0.222677)--(9.640098,0.17343)--(10.036205,0.117761)--cycle;
	\draw[color={black},line width = 0.5] (10.676402,4.672994)--(10.620732,4.276887)--(10.224625,4.332556)--(10.280294,4.728663)--cycle;
	\draw[color={black},line width = 0.5] (8.695866,4.95134)--(9.091973,4.895671)--(9.036303,4.499564)--(8.640196,4.555233)--cycle;
	\draw[color={black},line width = 0.5] (8,0)--(8.055669,0.396107)--(8.451776,0.340438)--(8.396107,-0.055669)--cycle;
	\draw [color={red}] (8.990268,-0.139173) node[anchor = center,scale=1, rotate = -8] {/};
\draw [color={red}] (9.686134,4.812167) node[anchor = center,scale=1, rotate = -8] {/};

	\draw [color={blue}] (10.328469,2.197324) node[anchor = center,scale=1, rotate = 82] {||};
\draw [color={blue}] (8.347933,2.47567) node[anchor = center,scale=1, rotate = -278] {||};

	\draw [color={black}] (8.920682,-0.634307) node[anchor = center,scale=1, rotate = -8.049059999999997] {2 cm};
	\draw [color={black}] (10.823603,2.127737) node[anchor = center,scale=1, rotate = -278.04906] {5 cm};

\end{tikzpicture}\\
}{6M2C}

 Calculer l'aire du rectangle.
\end{EXO}

% @see : https://coopmaths.fr/alea?uuid=c6fb2&id=6M2C&n=1&d=10&s=3&alea=RD7u&cd=1&cols=1
\begin{EXO}{\begin{tikzpicture}[baseline,scale = 0.75]

    \tikzset{
      point/.style={
        thick,
        draw,
        cross out,
        inner sep=0pt,
        minimum width=5pt,
        minimum height=5pt,
      },
    }
    \clip (-2,-3) rectangle (22,7);
    	\draw[color={black}] (15,0)--(17.963065,-0.469303)--(18.745237,4.469138)--cycle;
	\draw [color={black}] (14.573479,-0.254313) node[anchor = center,scale=1, rotate = 0] {N};
	\draw [color={black}] (18.14871,-0.929628) node[anchor = center,scale=1, rotate = 0] {O};
	\draw [color={black}] (18.962017,4.919588) node[anchor = center,scale=1, rotate = 0] {P};
	\draw[color={black},line width = 0.5] (17.963065,-0.469303)--(17.56799,-0.40673)--(17.630563,-0.011654)--(18.025639,-0.074228)--cycle;
	\draw [color={black}] (16.403315,-0.728496) node[anchor = center,scale=1, rotate = -9.022320000000008] {3 cm};
	\draw [color={black}] (18.847995,1.9217) node[anchor = center,scale=1, rotate = -279.08575] {5 cm};
	\draw [color={black}] (16.489393,2.555721) node[anchor = center,scale=1, rotate = 50.00583] {5,8 cm};

\end{tikzpicture}\\
}{6M2C}

 Calculer l'aire du triangle rectangle.
\end{EXO}

% @see : https://coopmaths.fr/alea?uuid=bae27&id=6M2C-flash1&n=1&d=10&alea=8VtO&cd=1&cols=1
\begin{EXO}{}{6M2C-flash1}

 Un carré de côté $8\text{ cm}$ et un rectangle de largeur $7\text{ cm}$ et de longueur $8\text{ cm}$ ont une aire qui diffère de $8\text{ cm}^2$. 
    
    	$\square\;$ Vrai\qquad $\square\;$ Faux\qquad \\

\end{EXO}

% @see : https://coopmaths.fr/alea?uuid=90e79&id=6M2C-flash2&n=1&d=10&alea=8N1g&cd=1&cols=1
\begin{EXO}{}{6M2C-flash2}

 Quelle est la longueur du côté d'un carré de $49\text{ cm}^2$ d'aire ?
\end{EXO}


\clearpage

\begin{Correction}
\begin{EXO}{}{}

 L'aire de cette figure est : ${\color[HTML]{F15929}\boldsymbol{9}}$ cm².
\end{EXO}

\begin{EXO}{}{}

 L'aire de cette figure est : ${\color[HTML]{F15929}\boldsymbol{4\times4\times4}}\text{ soit }{\color[HTML]{F15929}\boldsymbol{64}}\text{cm}^2$.
\end{EXO}

\begin{EXO}{}{}

 L'aire de la zone hachurée est $35$ cm² - $14$ cm² = ${\color[HTML]{F15929}\boldsymbol{21}}$ cm²
\end{EXO}

\begin{EXO}{}{}

 L'aire de la figure 1 est donnée par : ${\color[HTML]{F15929}\boldsymbol{6\times3\div2\times2}}\text{ soit }{\color[HTML]{F15929}\boldsymbol{18}}$ u.a, et l'aire de la figure 2 est donnée par : ${\color[HTML]{F15929}\boldsymbol{7\times4\div2\times2}}\text{ soit }{\color[HTML]{F15929}\boldsymbol{28}}$ u.a.
\end{EXO}

\begin{EXO}{}{}

 $1\text{ m}^2 = 100 ~\text{dm}^2$ et donc la mesure en $~\text{dm}^2$ est $100$ fois plus grande que la mesure en $\text{ m}^2$ ($1 \text{ m}^2$ est une centaine de $1 ~\text{dm}^2$).\\$0{,}05 \text{ m}^2 = 0{,}05 {\color[HTML]{216D9A}\boldsymbol{\times 1\,\textbf{ m}^2}}$ = $0{,}05 {\color[HTML]{216D9A}\boldsymbol{~\times~ 100\,~\textbf{dm}^2}} = {\color[HTML]{F15929}\boldsymbol{5}} ~\text{dm}^2$
\end{EXO}

\begin{EXO}{}{}

 $1~\text{dm}^2 = 100 \text{ cm}^2$ et donc la mesure en $\text{ cm}^2$ est $100$ fois plus grande que la mesure en $~\text{dm}^2$ ($1 ~\text{dm}^2$ est une centaine de $1 \text{ cm}^2$).\\$0{,}02 ~\text{dm}^2 = 0{,}02 {\color[HTML]{216D9A}\boldsymbol{\times 1\,~\textbf{dm}^2}}$ = $0{,}02 {\color[HTML]{216D9A}\boldsymbol{~\times~ 100\,\textbf{ cm}^2}} = {\color[HTML]{F15929}\boldsymbol{2}} \text{ cm}^2$
\end{EXO}

\begin{EXO}{}{}

 $1~\text{dm}^2 = \dfrac{1}{100} \text{ m}^2$ et donc la mesure en $\text{ m}^2$ est $100$ fois plus petite que la mesure en $~\text{dm}^2$ ($1 ~\text{dm}^2$ est un centième de $1 \text{ m}^2$).\\$0{,}04 ~\text{dm}^2 = 0{,}04 {\color[HTML]{216D9A}\boldsymbol{\times 1\,~\textbf{dm}^2}}$ = $0{,}04 {\color[HTML]{216D9A}\boldsymbol{\times \dfrac{1}{100}\,\textbf{ m}^2}}$ = $0{,}04 {\color[HTML]{216D9A}\boldsymbol{~\div~ 100\,\textbf{ m}^2}} = {\color[HTML]{F15929}\boldsymbol{0{,}000\,4}} \text{ m}^2$
\end{EXO}

\begin{EXO}{}{}

 $\mathcal{A}_{FGHI}=4\text{ cm}\times4\text{ cm}={\color[HTML]{F15929}\boldsymbol{16}}\text{ cm}^2$
\end{EXO}

\begin{EXO}{}{}

 $\mathcal{A}_{JKLM}=2\text{ cm}\times5\text{ cm}={\color[HTML]{F15929}\boldsymbol{10}}\text{ cm}^2$
\end{EXO}

\begin{EXO}{}{}

 $\mathcal{A}_{NOP}=3\text{ cm}\times5\text{ cm}\div2={\color[HTML]{F15929}\boldsymbol{7{,}5}}\text{ cm}^2$
\end{EXO}

\begin{EXO}{}{}

 {\bfseries \color[HTML]{f15929}Vrai}\\ $\bullet$  Le carré a une aire de $8\times 8=64\text{ cm}^2$.\\
    $\bullet$  Le rectangle a une aire de $8\times 7=56\text{ cm}^2$.\\Ce qui fait bien  une différence de $64 - 56=8\text{ cm}^2$ .
\end{EXO}

\begin{EXO}{}{}

 On sait que $49\text{ cm}^2 = 7\text{ cm} \times 7\text{ cm}$.\\Donc la longueur du côté du carré est de ${\color[HTML]{F15929}\boldsymbol{7}}\text{ cm}$.
\end{EXO}

\clearpage
\end{Correction}
\end{document}

% Local Variables:
% TeX-engine: luatex
% End: